% Chapter 4, Topic _Linear Algebra_ Jim Hefferon
%  http://joshua.smcvt.edu/linearalgebra
%  2012-Jun-18
\topic{Chiò 方法}
\index{Chi\`o's method|(}

对只含整数的矩阵做高斯消元时，人们常常希望保持这一特点。
为了在对下面这个矩阵化简时避免分数
\begin{equation*}
  A=
  \begin{mat}
    2 &1 &1 \\
    3 &4 &-1 \\
    1 &5 &1 
  \end{mat}
\end{equation*}
可以先把下面的行乘以~$2$
\begin{equation*}
  \grstep[2\rho_3]{2\rho_2}
  \begin{mat}
    2 &1 &1 \\
    6 &8 &-2 \\
    2 &10 &2 
  \end{mat}
  \tag{$*$}
\end{equation*}
使得第一列的消元过程如下进行。
\begin{equation*}
  \grstep[-\rho_1+\rho_3]{-3\rho_1+\rho_2}
  \begin{mat}
    2 &1 &1 \\
    0 &5 &-5 \\
    0 &8 &0 
  \end{mat}
  \tag{$**$}
\end{equation*}
这种全整数的做法便于心算。
而且，在计算机上用整数运算可以避免浮点计算带来的一些棘手问题 \cite{Kahan}。
因此这种做法有充分的理由。

这种做法的另一个优点是，我们可以方便地对~($**$) 的第一列用
Laplace 展开，然后只要记得除以~$4$（因为~($*$)）就能得到行列式。

下面是这种求行列式方法的一般 $\nbyn{3}$ 情形。
首先，假设 $a_{1,1}\neq 0$，我们可以把下面的行倍乘。
\begin{equation*}
  A=
  \begin{mat}
    a_{1,1} &a_{1,2} &a_{1,3}  \\
    a_{2,1} &a_{2,2} &a_{2,3}  \\
    a_{3,1} &a_{3,2} &a_{3,3}  
  \end{mat}
  \grstep[a_{1,1}\rho_3]{a_{1,1}\rho_2}
  \begin{mat}
    a_{1,1}       &a_{1,2}       &a_{1,3}        \\
    a_{2,1}a_{1,1} &a_{2,2}a_{1,1} &a_{2,3}a_{1,1}   \\
    a_{3,1}a_{1,1} &a_{3,2}a_{1,1} &a_{3,3}a_{1,1}   
  \end{mat}
\end{equation*}
这将行列式缩放了 $a_{1,1}^2$ 倍。
现在沿第一列消元。
\begin{equation*}
   \grstep[-a_{3,1}\rho_1+\rho_3]{-a_{2,1}\rho_1+\rho_2}
  \begin{mat}
    a_{1,1}       &a_{1,2}       &a_{1,3}       \\
    0 
       &a_{2,2}a_{1,1}-a_{2,1}a_{1,2} 
       &a_{2,3}a_{1,1}-a_{2,1}a_{1,3}               \\
    0 
       &a_{3,2}a_{1,1}-a_{3,1}a_{1,2} 
       &a_{3,3}a_{1,1}-a_{3,1}a_{1,3}                   
  \end{mat}
\end{equation*}
设 $C$ 是 $1,1$ 子式。
由 Laplace 公式，上述矩阵的行列式为 $a_{1,1}\det(C)$。
于是 $a_{1,1}^2\det(A)=a_{1,1}\det(C)$，又因 $a_{1,1}\neq 0$，
故得 $\det(A)=\det(C)/a_{1,1}$。

要处理更大的矩阵，我们必须弄清楚如何计算子式的元素。
上面的模式是：子式的每个元素都是一个
$\nbyn{2}$~行列式。
例如，子式左上角的元素 $a_{2,2}a_{1,1}-a_{2,1}a_{1,2}$，也就是上面矩阵中的 $2,2$~元素，
是 $A$ 的这四个
元素所成矩阵的行列式。
\begin{equation*}
  \begin{mat}
    \highlight{$a_{1,1}$} &\highlight{$a_{1,2}$} &a_{1,3}  \\
    \highlight{$a_{2,1}$} &\highlight{$a_{2,2}$} &a_{2,3}  \\
    a_{3,1}               &a_{3,2}               &a_{3,3}  
  \end{mat}
\end{equation*}
而子式的左下角，即上面的 $3,2$ 元素，
是矩阵 
中这四个元素
\begin{equation*}
  \begin{mat}
    \highlight{$a_{1,1}$} &\highlight{$a_{1,2}$} &a_{1,3}  \\
    a_{2,1}               &a_{2,2}               &a_{2,3}  \\
    \highlight{$a_{3,1}$} &\highlight{$a_{3,2}$} &a_{3,3}  
  \end{mat}
\end{equation*}

于是，当 $A$ 是 $n\geq 3$ 的~$\nbyn{n}$ 矩阵时，我们令 Chiò 矩阵 $C$ 为
$\nbyn{(n-1)}$ 矩阵，其 $i,j$ 元素是行列式
\begin{equation*}
  \begin{vmat}
    a_{1,1}  &a_{1,j+1} \\
    a_{i+1,1}    &a_{i+1,j+1}
  \end{vmat}
\end{equation*}
其中 $1<i,j\leq n$。
求 $A$ 的行列式的\definend{Chiò 方法}\index{Chi\`o's method} 
是：若 $a_{1,1}\neq 0$，则
$\det(A)=\det(C)/a_{1,1}^{n-2}$。
（顺便说一句，
Chiò 公式并不要求这些数是整数；它对实数同样适用。）

为了说明，我们求这个 $\nbyn{3}$ 矩阵的行列式。
\begin{equation*} 
  A=
  \begin{mat}
    2 &1 &1 \\
    3 &4 &-1 \\
    1 &5 &1 
  \end{mat}
\end{equation*}
这就是 Chiò 矩阵。
\begin{equation*}
  C=
  \begin{mat}
    \begin{vmat}
      2 &1 \\
      3 &4
    \end{vmat}
   &
   \begin{vmat}
     2 &1 \\
     3 &-1
   \end{vmat}        \\[3ex]      
   \begin{vmat}
     2 &1 \\
     1 &5
   \end{vmat}
   &
   \begin{vmat}
     2 &1 \\
     1 &1
   \end{vmat}
  \end{mat}
  =
  \begin{mat}
    5  &-5  \\
    9  &1
  \end{mat}
\end{equation*}
$\nbyn{3}$ 矩阵的公式
$det(A)=\det(C)/a_{1,1}$ 给出 $\det(A)=(50/2)=25$。

对于更大的行列式，我们必须做多个步骤，但每一步只涉及 $\nbyn{2}$ 行列式。
因此我们常常只需写下一点中间信息就能算出行列式。
例如，对这个 $\nbyn{4}$ 矩阵
\begin{equation*}
  A=
  \begin{mat}
    3  &0  &1  &1  \\
    1  &2  &0  &1  \\
    2  &-1 &0  &3  \\
    1  &0  &0  &1
  \end{mat} 
\end{equation*}
我们可以在心里完成每一个
$\nbyn{2}$ 计算，而只写下
$\nbyn{3}$ 的结果。
\begin{equation*}
  C_3=
  \begin{mat}
    \begin{vmat}
      3 &0 \\
      1 &2
    \end{vmat}
    &\begin{vmat}
     3 &1 \\
     1 &0 
    \end{vmat}
    &\begin{vmat}
     3 &1 \\
     1 &1
    \end{vmat}                \\[3ex]
    \begin{vmat}
     3 &0 \\
     2 &-1
    \end{vmat}
    &\begin{vmat}
     3 &1 \\
     2 &0
    \end{vmat}
    &\begin{vmat}
     3 &1 \\
     2 &3
    \end{vmat}             \\[3ex]
    \begin{vmat}
     3 &0 \\
     1 &0
    \end{vmat}
    &\begin{vmat}
     3 &1 \\
     1 &0
    \end{vmat}
    &\begin{vmat}
     3 &1 \\
     1 &1
    \end{vmat}
  \end{mat}
  =
  \begin{mat}
    6  &-1  &2 \\
   -3 &-2  &7 \\
    0  &-1  &2
  \end{mat}
\end{equation*}
注意这个矩阵的行列式
是~$A$ 的行列式的 $a_{1,1}^{4-2}=3^2$ 倍。

最后，迭代即可。
这是 $C_3$ 的 Chiò 矩阵。
\begin{equation*}
  C_2=
  \begin{mat}    
    \begin{vmat}
      6 &-1 \\
     -3 &-2 
    \end{vmat}
    &\begin{vmat}
      6 &2 \\
     -3 &7
    \end{vmat}           \\[3ex]
    \begin{vmat}
      6 &-1 \\
      0 &-1
    \end{vmat}
    &\begin{vmat}
      6 &2 \\
      0 &2
    \end{vmat}          
  \end{mat}
  =
  \begin{mat}
    -15 &48 \\
    -6 &12
  \end{mat}
\end{equation*}
这个矩阵的行列式
是~$C_3$ 的行列式的 $6$~倍。
$C_2$ 的行列式是 $108$。
于是
$\det(A)=108/(3^2\cdot 6)=2$。

Laplace 展开公式
把一个~$\nbyn{n}$ 行列式的计算化为若干个 $\nbyn{(n-1)}$~行列式的求值。
Chiò 公式也是递归的，
但它把一个~$\nbyn{n}$ 行列式化成单个~$\nbyn{(n-1)}$ 行列式，
而后者由若干个 $\nbyn{2}$ 行列式算得。
不过，对大矩阵而言，高斯消元法比这两种方法都好；
例如，
它所需的运算量大约是 Chiò 方法的一半~\cite{FullerLogan}。


\begin{exercises}
  \item 
    用 Chiò 方法求下列行列式。
    \begin{exparts*}
      \partsitem 
        $
        \begin{vmat}
          1  &2  &3  \\
          4  &5  &6  \\
          7  &8  &9  
        \end{vmat}
        $
      \partsitem 
        $
        \begin{vmat}
          2  &1  &4  &0  \\
          0  &1  &4  &0  \\
          1  &1  &1  &1  \\
          0  &2  &1  &1  
        \end{vmat}
        $
    \end{exparts*}
    \begin{answer}
      \begin{exparts*}
        \partsitem Chiò 矩阵是
        \begin{equation*}
          C=
          \begin{mat}
            -3 &-6 \\
            -6 &-12
          \end{mat}
        \end{equation*}
        其行列式为 $0$
        \partsitem 从
        \begin{equation*}
          C_3=
          \begin{mat}
            2 &8  &0 \\
            1 &-2 &2 \\
            4 &2  &2
          \end{mat}
        \end{equation*}
        开始，再下一步得
        \begin{equation*}
          C_2=
          \begin{mat}
            -12 &4 \\
            -28 &4
          \end{mat}
        \end{equation*}
        其行列式 $\det(C_2)=64$。 
        故原矩阵的行列式为 $64/(2^2\cdot 2^1)=8$
      \end{exparts*} 
    \end{answer}
  \item 如果 $a_{1,1}$ 是零怎么办？
    \begin{answer}
      与上面 $\nbyn{3}$~情形相同的构造表明，
      我们可以选取任一非零元素~$a_{i,j}$ 来代替 $a_{1,1}$。
      Chiò 矩阵的元素 $c_{p,q}$ 是行列式
      \begin{equation*}
        \begin{vmat}
          a_{1,1}    &a_{1,q+1} \\
          a_{p+1,1}  &a_{p+1,q+1}
        \end{vmat}
      \end{equation*}
      的值，其中 $p+1\neq i$ 且 $q+1\neq j$。
    \end{answer}
  \item \definend{Sarrus 法则}\index{Sarrus, Rule of}\index{Rule of Sarrus}
    是许多人学过的一个帮助记忆
    $\nbyn{3}$ 行列式公式的口诀。
    在矩阵右侧抄写前两列。
    \begin{equation*}
      \begin{array}{ccc|cc}
        a &b &c &a &b \\
        d &e &f &d &e \\
        g &h &i &g &h
      \end{array}
    \end{equation*}
    于是行列式就是三条左上到右下的对角线之和
    减去三条左下到右上的对角线之和：
    $aei+bfg+cdh-gec-hfa-idb$。
   请数一数 Sarrus 公式与 Chiò 方法各需要多少次运算。
   \begin{answer}
     Sarrus 公式用 $12$~次乘法和 $5$ 次加法（减法计入加法之中）。
     Chiò 公式对四个 $\nbyn{2}$~行列式中的每一个要用两次乘法和一次加法
     （实际上是减法），
     对 $\nbyn{2}$~Chiò 行列式还要再用两次乘法和一次加法，
     最后还要除以 $a_{1,1}$。
     总共是十一次乘除和五次加减。
     所以 Chiò 方法胜出。
   \end{answer}
  \item 证明 Chiò 公式。
    \begin{answer}
      考虑 $\nbyn{n}$ 矩阵
      \begin{equation*}
        A=
        \begin{mat}
           a_{1,1}  &a_{1,2}   &\cdots &a_{1,n-1}  &a_{1,n} \\
           a_{2,1}  &a_{2,2}   &\cdots &a_{2,n-1}  &a_{2,n} \\
                  &\vdots                         \\
           a_{n-1,1} &a_{n-1,2} &\cdots &a_{n-1,n-1} &a_{n-1,n}  \\ 
           a_{n,1}  &a_{n,2}   &\cdots &a_{n,n-1}  &a_{n,n} 
        \end{mat}
      \end{equation*}
      把除第一行外的每一行都乘以~$a_{1,1}$。
      \begin{equation*}
        \grstep[a_{1,1}\rho_3 \\ \vdots \\ a_{1,1}\rho_{n}]{a_{1,1}\rho_2}  
        \begin{mat}
          a_{1,1}         &a_{1,2}        &\cdots &a_{1,n-1}        &a_{1,n} \\
          a_{2,1}a_{1,1}   &a_{2,2}a_{1,1}  &\cdots &a_{2,n-1}a_{1,1}  &a_{2,n}a_{1,1} \\
                        &\vdots                         \\
          a_{n-1,1}a_{1,1} &a_{n-1,2}a_{1,1} &\cdots &a_{n-1,n-1}a_{1,1} &a_{n-1,n}a_{1,1}\\
          a_{n,1}a_{1,1}  &a_{n,2}a_{1,1}   &\cdots &a_{n,n-1}a_{1,1}  &a_{n,n}a_{1,1} 
        \end{mat}
      \end{equation*}
      这将行列式缩放了~$a_{1,1}^{n-1}$ 倍。 

      接下来对每个行~$i>1$ 执行行变换
      $-a_{i,1}\rho_1+\rho_i$。
      这些行变换不改变行列式。
      \begin{equation*}
        \grstep[-a_{3,1}\rho_1+\rho_3 \\ \vdotswithin{-a_{3,1}\rho_1+\rho_3} \\ -a_{n,1}\rho_1+\rho_n]{-a_{2,1}\rho_1+\rho_2}         
      \end{equation*}
      结果是一个矩阵~$B$，其第一行不变，第一列全为零（除了
      $1,1$ 元素为 $a_{1,1}$），其余元素如下。
      \begin{equation*}
        \begin{mat}
              a_{1,2}  
              &\cdots 
              &a_{1,n-1}  
              &a_{1,n} 
              \\
              a_{2,2}a_{1,1}-a_{2,1}a_{1,2}  
              &\cdots 
              &a_{2,n-1}a_{n,n}-a_{2,n-1}a_{1,n-1}  
              &a_{2,n}a_{n,n}-a_{2,n}a_{1,n}    
              \\
              \vdots                         \\
             a_{n,2}a_{1,1}-a_{n,1}a_{1,2} 
             &\cdots 
             &a_{n,n-1}a_{1,1}-a_{n,1}a_{1,n-1} 
             &a_{n,n}a_{1,1}-a_{n,1}a_{1,n}  
        \end{mat}
      \end{equation*}
       这个 $\nbyn{n}$ 矩阵 $B$ 的行列式是
       $A$ 的行列式的 $a_{1,1}^{n-1}$~倍。

     用~$C$ 表示矩阵 $B$ 的 $1,1$ 子式，
     即由最后 $n-1$ 行和最后 $n-1$ 列构成的子矩阵。
     沿 $B$ 的第一列作 Laplace 展开
     得到其行列式为 $\det(B) = (-1)^{1+1}a_{1,1}\det(C)$。

      若 $a_{1,1}\neq 0$，则令 $\set(B)$ 的两个表达式相等并约去公因子，
      得 $\det(A)=\det(C)/a_{n,n}^{n-2}$。       
    \end{answer}
\end{exercises}

\announcecomputercode
下面用计算机语言 Python 实现 Chiò 方法。
%  (to make it more readable 
% the code avoids some Python facilities).
% Note the recursive call in the final line of 
% \lstinline!chio_det!. 
\begin{lstlisting}
#!/usr/bin/python
# chio.py
#  Calculate a determinant using Chio's method.
# Jim Hefferon; Public Domain
# For demonstration only; for instance, does not handle the M[0][0]=0 case

def det_two(a,b,c,d):
    """Return the determinant of the 2x2 matrix [[a,b], [c,d]]"""
    return a*d-b*c

def chio_mat(M):
    """Return the Chio matrix as a list of the rows
        M  nxn matrix, list of rows"""
    dim=len(M)
    C=[]
    for row in range(1,dim):
        C.append([])
        for col in range(1,dim):  
            C[-1].append(det_two(M[0][0], M[0][col], M[row][0], M[row][col]))
    return C

def chio_det(M,show=None):
    """Find the determinant of M by Chio's method
        M  mxm matrix, list of rows"""
    dim=len(M)
    key_elet=M[0][0]
    if dim==1:
        return key_elet
    return chio_det(chio_mat(M))/(key_elet**(dim-2))


if __name__=='__main__':
    M=[[2,1,1], [3,4,-1], [1,5,1]]
    print "M=",M
    print "Det is", chio_det(M)
\end{lstlisting}
这是在命令行中调用该程序的结果。
\begin{lstlisting}
$ python chio.py
M=[[2, 1, 1], [3, 4, -1], [1, 5, 1]]
Det is 25
\end{lstlisting}
\index{Chi\`o's method|)}
\endinput
